\begin{abstract}
\thispagestyle{empty}

Выпускная квалификационная работа выполнена на тему: <<Веб-приложение для планирования питания с учетом предпочтений пользователей>>.

Назначение разработанного программного продукта состоит в автоматизации планирования рациона питания, формировании устойчивых здоровых пищевых привычек и повышении эффективности составления меню в повседневной практике пользователей.

Основные функции разработанного веб-приложения:
\begin{compactlist}
  \item автоматический расчет пищевой ценности для каждого рецепта;
  \item составление персональных планов питания на день, неделю или месяц (ручной и автоматический подбор);
  \item формирование списка покупок на основе выбранных рецептов;
  \item интеллектуальный ассистент для консультаций по питанию;
  \item аналитика пищевых привычек и прогресса достижения пользовательских целей.
\end{compactlist}

В работе используются следующие технологии: PostgreSQL (база данных), Python и Django (серверная часть), React и Effector (клиентская часть).

Работа включает введение, три главы, заключение, список использованных источников и приложения.

Во введении обоснована актуальность темы, сформулированы цель и задачи, определены объект и предмет исследования, а также практическая значимость результатов. В первой главе рассмотрена предметная область, проведен анализ целевой аудитории, ролевой модели и конкурентных решений, сформулированы требования к системе и обоснован выбор технологического стека. Во второй главе представлена реализация программного решения: проектирование структуры данных, разработка клиентской части, серверной части и прикладных модулей (планирование, расчет КБЖУ, модерация, уведомления). В третьей главе описаны тестирование, контейнеризация и порядок развертывания системы.

Практический результат ВКР --- работоспособное веб-приложение для планирования питания с поддержкой персонализации, автоматизации рутинных операций и ролевого управления контентом.

Практическая значимость разработанного решения состоит в снижении трудозатрат на составление рациона, повышении прозрачности контроля пищевой ценности и упрощении формирования списка покупок. Разработанная система может использоваться как самостоятельный цифровой инструмент для частных пользователей и как основа для дальнейшего развития сервисов в сфере здорового питания.
\end{abstract}